\subsection{Representation initialization: reserved endpoint}
\label{sec:pending-init}
\pendingslot{P01}{Complete the interrupted initialization experiment}{The
105-epoch record is partial; its observed maximum is not the 150-epoch endpoint.}

The interrupted arm initializes from the 300-epoch DeiT-S/8
\texttt{model\_best.pth.tar} file. Its startup log reports 152 loaded tensors,
25 missing tensors and no unexpected tensors. This is a partial mapping,
not evidence that every local and latent parameter inherits the dense model.
The checkpoint was selected using EMA accuracy, but the available loader
prefers \texttt{state\_dict}. The actual state key and tensor fingerprint
must be confirmed before calling this raw or EMA initialization. The local
encoder remains trainable; this is not router-only training.

The existing scratch reference with matching drop-path 0.07 is
\texttt{s5\_mn\_dp07\_150e}, whose best EMA top-1 is 79.534\%, subject to a
check of all other resolved arguments. The drop-path-0.1 selected arm is not
an initialization-only control. Report the source model's 300-epoch
pretraining separately from 150-epoch adaptation. Resume the interrupted
arm's own last checkpoint, retaining optimizer, scheduler, EMA and scaler
state; a fresh initialization would be a new run.

\begin{table}[htbp]\centering\small
\begin{tabular}{lrrrl}\toprule
Arm & Best EMA & Final EMA & Best epoch & State receipt\\\midrule
DeiT-initialized, completed 150 epochs & \pendingvalue & \pendingvalue & \pendingvalue & \pendingvalue\\
Difference to verified scratch & \pendingvalue & \pendingvalue & & \pendingvalue\\
\bottomrule\end{tabular}
\caption{Reserved initialization results [P01]. Accuracy is in percent and
differences in percentage points. Partial-curve values are not endpoints.}
\label{tab:pending-init}\end{table}
\pendingtext

\subsection{Progressive latent compression: reserved comparison}
\label{sec:pending-progressive}
\pendingslot{P02}{Test the queued schedule at a matched final budget}{The
configuration exists; no completed run or measured layerwise trace is present.}

The candidate retains the full-grid local encoder and the 784-to-392 gather.
It requests 192 additional removals across six latent blocks, with nominal
patch counts $392\to360\to328\to296\to264\to232\to200$.
The class token is additional. A runtime trace must confirm attention input,
merge output and MLP input counts. This candidate does not eliminate the
initial gather or shorten the earlier local blocks.

Comparing this arm with the native 392-patch model mixes schedule and final
budget. For schedule attribution, add a single-gather control at the same
achieved 200-patch endpoint, with matched depth, initialization, training
recipe and drop-path 0.07. Match ToMe's achieved budget for deployment
comparisons while recording its different training history.

\begin{table}[htbp]\centering\small
\begin{tabular}{lrrrr}\toprule
Arm & Final patches & Best EMA & Final EMA & Batch latency (ms)\\\midrule
Progressive latent & \pendingvalue & \pendingvalue & \pendingvalue & \pendingvalue\\
Single gather, same endpoint & \pendingvalue & \pendingvalue & \pendingvalue & \pendingvalue\\
ToMe, same endpoint & \pendingvalue & \pendingvalue & \pendingvalue & \pendingvalue\\
\bottomrule\end{tabular}
\caption{Reserved schedule comparison [P02]. Requested budgets do not
substitute for achieved counts. Timing requires the common protocol in P03.}
\label{tab:pending-progressive}\end{table}
\pendingtext
